\documentclass{report}
\usepackage{amsmath,amssymb}
\setlength{\parindent}{0mm}
\setlength{\parskip}{1em}
\begin{document}
\begin{center}
\rule{6in}{1pt} \
{\large Hormozd Gahvari \\
{\bf Algebraic Multigrid Solvers on Future Supercomputers}}

Department of Computer Science \\ University of Illinois at Urbana-Champaign \\ 201 N Goodwin Avenue \\ Urbana \\ IL 61801
\\
{\tt gahvari@illinois.edu}\\
William Gropp\\
Kirk E. Jordan\\
Luke Olson\\
Martin Schulz\\
Ulrike Meier Yang\end{center}

With single-core speeds no longer rising, the large supercomputers of the
future will feature increasing levels of parallelism. Applications will
have to deal with this reality in order to scale effectively to these
machines. In this talk, we focus on how to scale algebraic multigrid to
future machines. The big challenge is an increasing amount of
communication on coarse grids. We overview the body of past work on the
subject, and present performance models we are using to guide our work to
avoid difficulties that the past work faced. We then present results for
two approaches, using hybrid distributed memory/shared memory programming
and redundantly distributing data across otherwise inactive processors on
coarse grids.


\end{document}
